% META: {"id": "1NSI-ALGO-PARCOURS-TRIS-COURS-04", "chapitre": "1NSI-ALGO-PARCOURS-TRIS", "type_objet": "cours", "capacites": ["1NSI-ALGO-PARCOURS-TRIS-C4"], "mode_creation": "ex_nihilo", "sources_inspiration": [], "status": "approved", "capacites_codes": ["C4"], "statut": "structure"}
\section{Recherche dichotomique}

\margeAppui{Chercher un mot dans un dictionnaire papier trié : on ouvre au milieu, on
compare, et on ne garde que la moitié pertinente. C'est le principe de la
\textbf{dichotomie}.}

\definition[D1]{
  La \textbf{recherche dichotomique} recherche une valeur dans un tableau \textbf{trié}
  en comparant systématiquement la valeur cherchée à l'élément \textbf{du milieu} de la
  zone de recherche restante, et en ne conservant que la moitié où la valeur peut encore se
  trouver.
}

\begin{codereference}{Recherche dichotomique}
\begin{python}
def recherche_dichotomique(tableau_trie, valeur):
    """Renvoie l'indice de `valeur` dans `tableau_trie` (trie), ou -1 si absente."""
    borne_inf, borne_sup = 0, len(tableau_trie) - 1
    while borne_inf <= borne_sup:
        milieu = (borne_inf + borne_sup) // 2
        if tableau_trie[milieu] == valeur:
            return milieu
        elif tableau_trie[milieu] < valeur:
            borne_inf = milieu + 1
        else:
            borne_sup = milieu - 1
    return -1

t = [1, 3, 5, 7, 9, 11, 13, 15]
print(recherche_dichotomique(t, 11))  # 5
print(recherche_dichotomique(t, 4))   # -1
\end{python}
\end{codereference}

% BEGIN-VERIFY
% def recherche_dichotomique(tableau_trie, valeur):
%     borne_inf, borne_sup = 0, len(tableau_trie) - 1
%     while borne_inf <= borne_sup:
%         milieu = (borne_inf + borne_sup) // 2
%         if tableau_trie[milieu] == valeur:
%             return milieu
%         elif tableau_trie[milieu] < valeur:
%             borne_inf = milieu + 1
%         else:
%             borne_sup = milieu - 1
%     return -1
% t = [1, 3, 5, 7, 9, 11, 13, 15]
% assert recherche_dichotomique(t, 11) == 5
% assert recherche_dichotomique(t, 4) == -1
% assert recherche_dichotomique(t, 1) == 0
% assert recherche_dichotomique(t, 15) == 7
% END-VERIFY

\propriete[Coût de la dichotomie]{
  À chaque itération, la zone de recherche est \textbf{divisée par deux}. Le coût de la
  recherche dichotomique est donc \textbf{logarithmique} ($O(\log_2 n)$), bien plus rapide
  que le coût linéaire de la recherche séquentielle pour un grand tableau.
}

\demonstration{
\textbf{Terminaison de la recherche dichotomique.}

On considère le \textbf{variant de boucle} $V=\text{borne\_sup}-\text{borne\_inf}$.

\textbf{$V$ est bien défini et positif ou nul} tant que la boucle continue, puisque la
condition de la boucle \lstinline{while} est précisément \lstinline{borne_inf <= borne_sup}.

\textbf{$V$ décroît strictement} à chaque itération : si \lstinline{tableau_trie[milieu] <
valeur}, alors \lstinline{borne_inf} devient \lstinline{milieu + 1}, strictement supérieur à
son ancienne valeur (puisque $\text{milieu}\geqslant\text{borne\_inf}$), donc $V$ diminue.
Si \lstinline{tableau_trie[milieu] > valeur}, alors \lstinline{borne_sup} devient
\lstinline{milieu - 1}, strictement inférieur à son ancienne valeur (puisque
$\text{milieu}\leqslant\text{borne\_sup}$), donc $V$ diminue également.

\textbf{Conclusion.} $V$ est un entier positif ou nul qui décroît strictement à chaque
tour de boucle : une suite d'entiers positifs strictement décroissante est nécessairement
\textbf{finie}. La boucle \lstinline{while} termine donc en un nombre fini d'itérations.
}

% BEGIN-VERIFY
% def recherche_dichotomique_avec_variant(tableau_trie, valeur):
%     borne_inf, borne_sup = 0, len(tableau_trie) - 1
%     variants = []
%     while borne_inf <= borne_sup:
%         variants.append(borne_sup - borne_inf)
%         milieu = (borne_inf + borne_sup) // 2
%         if tableau_trie[milieu] == valeur:
%             return milieu, variants
%         elif tableau_trie[milieu] < valeur:
%             borne_inf = milieu + 1
%         else:
%             borne_sup = milieu - 1
%     return -1, variants
% t = [1, 3, 5, 7, 9, 11, 13, 15]
% _, variants = recherche_dichotomique_avec_variant(t, 4)
% for k in range(1, len(variants)):
%     assert variants[k] < variants[k - 1]
% assert all(v >= 0 for v in variants)
% END-VERIFY

\erreurFrequente{Appliquer la dichotomie sur un tableau \textbf{non trié}. La dichotomie
repose entièrement sur le tri : comparer à l'élément du milieu ne permet d'éliminer une
moitié \textbf{que si} on est certain que tous les éléments de cette moitié sont, eux
aussi, inférieurs (ou supérieurs) à la valeur cherchée — ce qui n'est garanti que si le
tableau est trié.}

\alamachine{Modifie \lstinline{recherche_dichotomique} pour qu'elle compte le nombre
d'itérations effectuées, et compare ce nombre à $\log_2(n)$ pour un tableau trié de $1000$
éléments (utilise \lstinline{math.log2}).}
